\section{Perverse cohomology of a dense eigencomplex}
\label{sec:perverse}

In this section the curve is defined over $\mathbb C$.  We prove that
perverse cohomology preserves the entire Hecke eigenstructure for a
Zariski-dense parameter.  The main point is to obtain exactness near the
given parameter from a frame on its free closed orbit.  The spectral action
and local constancy needed to use this frame are the Betti results of
Nadler--Yun.

Put $N_B=R_u(B)$ and $T=B/N_B$.  Write
\[
 \cA=\Bun_{G,B,x}(X),\qquad
 \cA^+=\Bun_{G,N_B,x}(X).
\]
Thus $\cA^+$ remembers an enhancement of the flag, namely an
$N_B$-reduction.  Let $\cB$ denote either $\cA$ or $\cA^+$, and let
$\Lambda\subset T^*\cB$ be its cone of generically nilpotent Higgs fields,
with the residue condition imposed by the chosen level.  All Hecke functors
below act away from $x$ and use the relative normalization of the
introduction.

\begin{proposition}\label{prop:perverse-eigen}
Let $X/\mathbb C$ be a smooth projective connected curve, let $x\in X$,
and put $U=X\setminus\{x\}$.  Let $G$ be simple and simply connected,
and let $\cB$ be either of the two level stacks just defined.  Suppose
that a locally bounded constructible geometric $E$-adic complex $\cF$
on $\cB$ has a coherent Hecke eigenstructure for a Zariski-dense
parameter $\sigma$ on $U$.  For every $j\in\mathbb Z$, the perverse
cohomology ${}^pH^j(\cF)$ has an induced coherent eigenstructure with
the same parameter and has singular support in $\Lambda$.  This
structure is natural in representations and compatible with the unit,
convolution, permutations, and all collision diagonals in every $U^I$.
If $\cF\ne0$, at least one of these perverse eigensheaves is nonzero.
No common bound on the cohomological degrees of $\cF$ over $\cB$ is
required.
\end{proposition}

\subsection{Betti transport and the spectral action}

We first record precisely which Betti statements enter the proof.
For a formal coordinate $t$ at $x$, the ordinary and enhanced levels
correspond respectively to
\[
 \mathrm{Iw}=\{g\in G(\mathbb C[[t]]):g(0)\in B\},\qquad
 \mathrm{Iw}^{0}=\{g\in G(\mathbb C[[t]]):g(0)\in N_B\}.
\]
Both contain the first congruence subgroup and are contained in the
positive-loop maximal parahoric, so they satisfy the level hypotheses
of \cite[\S6.2]{NY}.  The nilpotent cone in that section is defined
by generic nilpotence of the Higgs field in the cotangent stack with
level; it is therefore our $\Lambda$.

Let $\mathscr D_\Lambda(\cB)$ be the Betti category of complexes with
singular support in $\Lambda$.  Nadler--Yun's
\cite[Theorem~6.2.2]{NY} gives, for every finite set $I$ and collection
$\boldsymbol V=(V_i)_{i\in I}$,
\begin{equation}\label{eq:perverse-NY-support}
 \SS\bigl(\Hecke_{I,\boldsymbol V}(K)\bigr)
       \subset \Lambda\times 0_{U^I},
       \qquad K\in\mathscr D_\Lambda(\cB).
\end{equation}
This assertion holds on the full product $U^I$, including its diagonals.
Their product-stratification argument
\cite[Proposition~6.3.2]{NY} consequently gives locally constant
transport in the leg variables, compatibly with the multi-point tensor
operations.  Explicitly, on a smooth finite-type chart $D$ of $\cB$ one
chooses a stratification whose conormals contain the pulled-back
nilpotent cone.  On the product with a sufficiently small contractible
open set $P\subset (U^{\mathrm{an}})^I$, the output is weakly
constructible for the product stratification and is the pullback of its restriction
to any slice $D\times\{\boldsymbol u\}$.  The same statement applies
when $P$ meets collision diagonals.

Choose $u\in U^{\mathrm{an}}$ and free generators
$\gamma_1,\ldots,\gamma_a$ of
$\pi_1(U^{\mathrm{an}},u)$; here $a=2g(X)$ because there is one
puncture.  Set
\begin{equation}\label{eq:perverse-spectral-stack}
 Y=\Gd^{a},\qquad
 \mathscr L=[Y/\Gd],
\end{equation}
where $\Gd$ acts by simultaneous conjugation.  Since the punctured
surface has the homotopy type of a graph, $\mathscr L$ is its Betti
stack of $\Gd$-local systems.  By
\cite[Theorem~6.3.7]{NY}, the symmetric monoidal category
$\Perf(\mathscr L)$ acts on $\mathscr D_\Lambda(\cB)$.  For
$V\in\Rep(\Gd)$, the equivariant vector bundle
\[
 \mathscr E_V=Y\times V
\]
acts by the fixed-point Hecke functor $\Hecke_{V,u}$.  Its
tautological automorphisms at the chosen generators act by Hecke
transport along $\gamma_i$.  We identify Satake and its eigenvalues
with these conventions, as in Section~\ref{sec:foundations}.
These theorems allow all complexes in the indicated nilpotent category:
the absence of global boundedness restrictions is explicit in
\cite[\S5.1 and \S6.2]{NY}.  In particular their application here does
not require support in a finite-type substack.  Our finite-dimensional
Satake kernels also preserve local bounded constructibility: on a
finite-type output chart, each bounded Hecke correspondence is proper
of finite type and meets only a quasi-compact part of the input stack.

A trivialization of a parameter $\sigma$ at $u$ represents its
topological monodromy by a homomorphism
$\rho:\pi_1(U^{\mathrm{an}},u)\to\Gd(E)$.  Its values on the
chosen generators give the corresponding point $y\in Y(E)$.

\begin{lemma}\label{lem:perverse-scalar-center}
Let $K\in\mathscr D_\Lambda(\cB)$ be a coherent Betti Hecke
eigencomplex with parameter $\sigma$, represented by
$y=(\rho(\gamma_1),\ldots,\rho(\gamma_a))\in Y(E)$.
The degree-zero central endomorphism associated to
$f\in E[Y]^{\Gd}=\operatorname{End}^0_{\Perf(\mathscr L)}(\cO)$
acts on $K$ as $f(y)\operatorname{id}_K$.
\end{lemma}

\begin{proof}
The identification of spectral functions with excursion operators is
\cite[\S7.1 and Proposition~7.2.3]{NY}.  In the present free-group
case its evaluation rule can be seen directly.  Matrix coefficients
span the coordinate ring of $Y$.  Exactness of invariants for the
reductive group $\Gd$ shows that an invariant function is a finite sum
of functions of the form
\begin{equation}\label{eq:perverse-excursion-trace}
 (g_1,\ldots,g_a)\longmapsto
 \operatorname{Tr}_{W}
 \left(A\circ\bigotimes_{i=1}^a V_i(g_i)\right),
 \qquad
 W=\bigotimes_{i=1}^a V_i,\quad A\in\operatorname{End}_{\Gd}(W).
\end{equation}
Indeed $E[Y]=\bigotimes_i E[\Gd]$ is spanned by products of
matrix coefficients.  Each such product lies in the image of the
equivariant map $\operatorname{End}_E(W)\to E[Y]$ given by the
trace in \eqref{eq:perverse-excursion-trace}.  A given invariant
function belongs to the image of finitely many of these maps.
The direct sum of their source spaces surjects onto a
finite-dimensional $\Gd$-stable image containing that function.
Exactness of invariants lifts it to a tuple of invariant
endomorphisms, since
$(\bigoplus_\nu\operatorname{End}_E(W_\nu))^{\Gd}
=\bigoplus_\nu\operatorname{End}_{\Gd}(W_\nu)$.
This proves the asserted finite-sum expression.

The corresponding excursion inserts the invariant coevaluation tensor
for $W\otimes W^*$, transports the $V_i$ factors around $\gamma_i$,
applies $A$, and evaluates.  The $W^*$ factor may be retained at the
base point as an additional identity leg.  Naturality, convolution,
and the unit constraints of the eigenstructure identify these maps
with the same operations on the tensor local system of $\sigma$.
Their composite on the multiplicity space is precisely
\eqref{eq:perverse-excursion-trace} evaluated at $y$.  Summing proves
the assertion.  This uses the maps in the coherent eigenstructure,
not just the isomorphism classes of its fixed-point transforms.
\end{proof}

\subsection{A frame near the dense parameter}

The next lemma gives an invariant function nonzero at $y$ whose inversion
trivializes $\mathscr E_V$.  Through the spectral action, this will identify
$\Hecke_{V,u}$ with $\dim V$ copies of the identity on a subcategory
preserved by perverse truncation.  We only need the frames to prove this
exactness; they need not be compatible with tensor products.  The induced
eigenmaps will instead come from truncating the original coherent maps.

\begin{lemma}\label{lem:perverse-frame}
Suppose that $y\in Y(E)$ generates a Zariski-dense subgroup of
$\Gd$.  For every nonzero $V\in\Rep(\Gd)$ of dimension $b$ there
are $f_V\in E[Y]^{\Gd}$ with $f_V(y)\ne0$ and equivariant maps
\[
 \alpha_V:\cO_Y^{\oplus b}\longrightarrow\mathscr E_V,
 \qquad
 \beta_V:\mathscr E_V\longrightarrow\cO_Y^{\oplus b}
\]
such that
\begin{equation}\label{eq:perverse-adjugate}
 \beta_V\alpha_V=f_V\operatorname{id}_{\cO_Y^{\oplus b}},
 \qquad
 \alpha_V\beta_V=f_V\operatorname{id}_{\mathscr E_V}.
\end{equation}
\end{lemma}

\begin{proof}
The group $\Gd$ is adjoint because $G$ is simply connected.  The
stabilizer of $y$ centralizes a dense subgroup and hence equals
$Z(\Gd)=1$.  Its orbit $O$ is closed.  To recall the closed-orbit
criterion in this case, a one-parameter subgroup producing a limit
under simultaneous conjugation forces every entry of the tuple to
belong to its associated parabolic.  Density forces that parabolic
to be $\Gd$, so the one-parameter subgroup is central.  The
Hilbert--Mumford criterion therefore gives closedness.  Thus
$O\simeq\Gd$ and the equivariant bundle $\mathscr E_V|_O$ has
equivariant sections $s_1,\ldots,s_b$ specified by any basis
$v_1,\ldots,v_b$ of $V$ at $y$: at $g\cdot y$ their values are
$g v_1,\ldots,g v_b$.

Because $O$ is closed in the affine variety $Y$, restriction gives a
surjection $E[Y]\otimes V\to E[O]\otimes V$.  Taking invariants
remains exact in characteristic zero.  The $s_i$ therefore extend to
equivariant regular sections of $\mathscr E_V$ on $Y$.  Let
$\alpha_V$ have these sections as its columns.  Every algebraic
character of the semisimple group $\Gd$ is trivial, so
$\det(V)$ is trivial.  After fixing a volume form, the determinant
$f_V=\det(\alpha_V)$ is an invariant function, nonzero at $y$.
The adjugate matrix defines $\beta_V$ regularly on all of $Y$.
It is equivariant, either by the exterior-power construction of the
adjugate or by the identity
$\beta_V=f_V\alpha_V^{-1}$ on the open set where $f_V\ne0$.
The adjugate identities give \eqref{eq:perverse-adjugate}.
\end{proof}

We now use these algebraic maps to obtain exactness on a subcategory
containing the eigencomplex and all of its truncations.  Write
$\mathscr D^{\mathrm{lcc}}_\Lambda(\cB)$ for the locally bounded
constructible objects of the Betti nilpotent category.  Perverse
truncation preserves this category: locally on a smooth chart,
singular support is unchanged upon replacing a complex by the union
of the singular supports of its perverse cohomologies.  This is the
usual microlocal d\'evissage, or equivalently the exactness of the
perverse vanishing-cycle tests.

\begin{lemma}\label{lem:perverse-localized-exactness}
Choose a frame as in Lemma~\ref{lem:perverse-frame} for each
isomorphism class of nonzero representations, and let $S$ be the
multiplicative set generated by the $f_V$.  The full subcategory
\[
 \mathscr D_S=
 \{K\in\mathscr D^{\mathrm{lcc}}_\Lambda(\cB):
        f_K:K\longrightarrow K\text{ is invertible for all }f\in S\}
\]
is preserved by perverse truncations and by all fixed-point Hecke
functors.  Every fixed-point Hecke functor is perverse $t$-exact on
$\mathscr D_S$.
\end{lemma}

\begin{proof}
For any central degree-zero natural endomorphism $f$ of the identity,
naturality for ${}^p\tau_{\le n}K\to K$, together with the universal
property of truncation, gives
\begin{equation}\label{eq:perverse-center-truncation}
 f_{{}^p\tau_{\le n}K}={}^p\tau_{\le n}(f_K).
\end{equation}
For example the map on the left is the unique endomorphism of
${}^p\tau_{\le n}K$ whose composite with its map to $K$ is
$f_K$ composed with that map.  The same argument for the other
truncation gives
$f_{{}^p\tau_{\ge n}K}={}^p\tau_{\ge n}(f_K)$.
Consequently invertibility of $f_K$ implies invertibility on both
truncations.  This establishes the asserted restricted
$t$-structure without any exactness assertion for arbitrary spectral
operators.

Apply the spectral action to the two maps of
Lemma~\ref{lem:perverse-frame}.  They give natural maps
\[
 K^{\oplus b}\xrightarrow{\alpha_{V,K}}
 \Hecke_{V,u}(K)
 \xrightarrow{\beta_{V,K}}K^{\oplus b}.
\]
The composites are the indicated actions of $f_V$.  Symmetric
monoidality identifies the action of any $f\in S$ on
$\Hecke_{V,u}(K)$ with $\Hecke_{V,u}(f_K)$.
Thus $\Hecke_{V,u}$ preserves $\mathscr D_S$, and
$f_V^{-1}\beta_{V,K}$ is an inverse to $\alpha_{V,K}$ there.
In particular
\begin{equation}\label{eq:perverse-fixed-frame}
 \Hecke_{V,u}|_{\mathscr D_S}
         \simeq \operatorname{id}_{\mathscr D_S}^{\oplus b}.
\end{equation}
This proves $t$-exactness at $u$.  Locally constant Hecke transport
along a path from $u$ to any other point gives the same conclusion
there.  The zero representation acts by zero.  Compositions of
fixed-point Hecke functors, including convolutions at one point,
are consequently $t$-exact and preserve $\mathscr D_S$.
\end{proof}

\subsection{Truncation of the moving Hecke system}

Fixed-point exactness now supplies the perverse bounds.  Locally constant
transport will extend them over every leg space, including its collision
diagonals.  We then use the resulting truncation comparisons to transport
the original eigenmaps and their coherence constraints.

\begin{proof}[Proof of Proposition~\ref{prop:perverse-eigen}]
Pass temporarily to Betti realization, with coefficients in $E$.
Theorem~\ref{thm:nilpotent-detection} puts $\cF$ in the nilpotent
category.  Its parameter remains Zariski dense after restriction
from the profinite fundamental group to topological loops.  Indeed
the latter have dense image in the former; any polynomial equation
on the matrices is closed for the adic topology after passing to a
finite coefficient extension containing its coefficients.  An
equation vanishing on topological monodromy therefore vanishes on
the entire adic image.

Use this tuple $y$ in Lemma~\ref{lem:perverse-frame}.
Lemma~\ref{lem:perverse-scalar-center} gives
$f_{V,\cF}=f_V(y)\operatorname{id}_{\cF}$, so $\cF\in\mathscr D_S$.
All its perverse truncations and cohomologies belong to this same
subcategory by Lemma~\ref{lem:perverse-localized-exactness}.
We next extend fixed-point exactness to the moving system and
explain its compatibility with collisions.

For $K\in\mathscr D_S$ and a finite set $I$, put $m=|I|$ and consider
\[
 \mathrm T_{I,\boldsymbol V}(K)
       =\Hecke_{I,\boldsymbol V}(K)[m].
\]
On a contractible parameter patch as in
\eqref{eq:perverse-NY-support}, its unshifted value is the exterior
product of a fixed-point transform and the constant sheaf of the
patch.  That transform is a succession of fixed-point Hecke
operators; at a collision it is their convolution.  By
Lemma~\ref{lem:perverse-localized-exactness} it has the same
perverse bounds as $K$.  The constant sheaf shifted by $m$ is
perverse on the smooth complex $m$-fold parameter space.  Testing
on smooth charts therefore proves that
$\mathrm T_{I,\boldsymbol V}$ is $t$-exact on $\mathscr D_S$.
In particular, the images of a truncation triangle for $K$ are
the truncation triangle for $\mathrm T_{I,\boldsymbol V}(K)$.
Its universal property supplies natural comparison isomorphisms
\begin{equation}\label{eq:perverse-moving-truncation}
 {}^pH^j\bigl(\Hecke_{I,\boldsymbol V}(K)[m]\bigr)
     \simeq \Hecke_{I,\boldsymbol V}({}^pH^jK)[m].
\end{equation}

These comparisons also hold for successive Hecke operations when
the input already carries other leg variables.  Locally in all
the parameter variables it is a constant family, so the same
fixed-point exactness argument applies, with the shift equal to
the total dimension of the retained smooth parameter space.
This proves the comparison for convolution diagrams as well as
for individual functors.

More explicitly, for a surjection $\pi:I\twoheadrightarrow J$ let
$\Delta_\pi:U^J\to U^I$ be the collision diagonal and put
$W_j=\bigotimes_{i\in\pi^{-1}(j)}V_i$.  Fusion gives the ordinary
pullback identification
\begin{equation}\label{eq:perverse-diagonal-normalization}
 (\operatorname{id}\times\Delta_\pi)^*
       \Hecke_{I,\boldsymbol V}(K)
       \simeq \Hecke_{J,\boldsymbol W}(K).
\end{equation}
On the locally constant families under consideration, the functor
\[
 (\operatorname{id}\times\Delta_\pi)^*[\,|J|-|I|\,]
\]
is perverse $t$-exact: in product charts it replaces the perverse
constant sheaf on an $|I|$-dimensional smooth base by the perverse
constant sheaf on an $|J|$-dimensional smooth base.  Applied to
\eqref{eq:perverse-diagonal-normalization}, this proves compatibility
of \eqref{eq:perverse-moving-truncation} with diagonal restriction.
The full-product assertion in \eqref{eq:perverse-NY-support} is
essential here; constancy only at pairwise distinct legs would
not provide this comparison.

Now let
$\mathcal V_{I,\sigma}=\boxtimes_{i\in I}(V_i)_\sigma$.
Shift the original eigenisomorphism by $m$ and apply ${}^pH^j$.
Exterior product with the lisse sheaf
$\mathcal V_{I,\sigma}[m]$ is $t$-exact, so
\eqref{eq:perverse-moving-truncation} gives
\[
 \Hecke_{I,\boldsymbol V}({}^pH^j\cF)[m]
   \simeq
 ({}^pH^j\cF\boxtimes\mathcal V_{I,\sigma})[m].
\]
Canceling the common shift yields precisely the relative-normalized
eigenisomorphism
\begin{equation}\label{eq:perverse-induced-eigen}
 \Hecke_{I,\boldsymbol V}({}^pH^j\cF)
   \simeq {}^pH^j\cF\boxtimes\mathcal V_{I,\sigma}.
\end{equation}
There is no moving-leg shift in this formula.

To verify coherence, apply these functorial truncation comparisons
to each diagram for the original eigenstructure.  Naturality in
representations, the unit, and permutations commute with the
comparisons by their construction from truncation.  The relative
exactness for successive operations gives convolution compatibility;
the dimension-adjusted diagonal comparison gives every fusion
compatibility, including iterated collisions.  On a diagram over
$U^I$, shift each entire vertex by $|I|$ once: successive
modifications at a shared leg do not introduce further leg
parameters.  Every vertex is locally the exterior product of a
composite of fixed-point exact functors applied to ${}^pH^j\cF$
with a lisse multiplicity space.  Thus all these shifted vertices
lie in the perverse heart.  After a collision to $U^J$, the shift
$|J|-|I|$ puts all vertices in the corresponding target heart.
For two objects $P,Q$ of any of these hearts,
$\pi_r\operatorname{Map}(P,Q)=\operatorname{Hom}(P,Q[-r])=0$
for $r>0$.  Their mapping spaces are therefore discrete; positive
Ext groups do not affect this assertion.  The resulting
commutative diagrams consequently determine the required coherent
constraints.  In particular choices of the auxiliary spectral
frames introduce no choices into \eqref{eq:perverse-induced-eigen}.

Finally these constructions give eigenmaps in geometric adic
sheaves, not only after realization.  Start with the adic
truncation triangles and apply the adic Hecke functors.  Betti
realization, which preserves perversity and detects isomorphisms
on finite-type charts, verifies the perverse bounds just proved.
The universal property of truncation then gives
\eqref{eq:perverse-moving-truncation} in the adic category itself.
The same reasoning applies to successive operations and diagonal
pullbacks.  Applying it to the original adic eigenmaps therefore
constructs \eqref{eq:perverse-induced-eigen} and all its
compatibilities there.  Singular support remains in $\Lambda$
by the perverse d\'evissage already used, or by
Theorem~\ref{thm:nilpotent-detection} applied to these eigenmaps.

Every assertion has been checked on finite-type smooth charts and
is compatible with further smooth pullback.  Thus no uniform
cohomological bound on the bundle stack has entered the proof.
If $\cF$ is nonzero, its restriction to some such chart is
nonzero and bounded.  Some perverse cohomology on that chart is
nonzero, and smooth perverse pullback identifies it with a
restriction of ${}^pH^j\cF$.  This proves the final assertion.
\end{proof}
